\begin{titlepage}
\centering
{\Large 深度访谈教学笔记\par}
\vspace{1.0cm}
{\huge\bfseries Gemini Robotics、跨本体与机器人世界模型\par}
\vspace{0.5cm}
{\Large Zhang Xiaojun 对话 Google DeepMind 谭杰\par}
\vspace{0.35cm}
{\large 基于 Zhang Xiaojun Podcast 公开视频、GPU 转写稿、章节结构与自制概念图整理\par}
\vspace{0.3cm}
{\large 2025-11-28\par}
\vspace{0.85cm}
\includegraphics[width=0.88\textwidth,height=0.42\textheight,keepaspectratio]{cover.jpg}\par
\vfill
\begin{tcolorbox}[width=0.92\textwidth, colback=black!2!white, colframe=black!55, sharp corners]
\textbf{视频标题}: 121. 对DeepMind谭捷的访谈: 机器人、跨本体、世界模型、Gemini Robotics 1.5和Google\par
\textbf{视频频道}: Zhang Xiaojun Podcast\par
\textbf{Duración del video}: 02:06:16\par
\textbf{Enlace del video}: \href{\videourl}{\nolinkurl{\videourl}}\par
\textbf{素材说明}: YouTube 无可用字幕；正文基于 faster-whisper large-v3 CUDA 转写、视频章节、封面和自制教学图整理.固定访谈画面不在正文重复使用.\par
\textbf{Página del proyecto}: \href{\repourl}{\nolinkurl{\repourl}}\par
\end{tcolorbox}
\end{titlepage}

\tableofcontents
\newpage

